\section{What a linear-time transposer would need}\label{sec:construction}

Our aim is to decide~(T), and we would prefer a construction. This section collects what the results above require of
any linear-time transposer, and of any oblivious one in particular. It contains no new theorems: every claim is one of
the results above or the first two papers, a citation, or explicitly labelled as a heuristic.

\subsection{Necessary conditions}

A machine witnessing~(T) must have all of the following properties, in addition to those of the first two
papers~\cite{FirstPaper2026,SecondPaper2026} ($\Omega(\log d)$ passes, folding macroscopic motion, and a coding gain of order
$\log d$ over routing on the block network of a slightly modified machine; within the abstract local-coding models of
the second paper, a fixed fractional saving is also excluded when the interaction graph has treewidth $o(d)$).
\begin{enumerate}[label=(C\arabic*),leftmargin=3em]
\item \emph{Curved motion.} In every macroscopic limit, the work heads travel a total distance of at least
half a block unit, a linear number of cells, along motion that is not countably-piecewise affine (Theorem~Cq,
\cref{thm:Cq}). In every clock, that is, every reparametrization of time by a homeomorphism whose inverse preserves null
sets, and for every countable set of affine maps, at least one block unit of travel has first returns that are not given
by these maps (Theorem~Cr, \cref{thm:Cr}); so neither a common change of clock nor fixed changes of scale remove the
curvature (\cref{cor:mot-clock}). The curved motion must be head-dependent, and
at least two work-head limits must be non-monotone and must not all have the same level sets (\cref{cor:mot-design}). Curved monotone expanders of the Bourgain--Yehudayoff type~\cite{BourgainYehudayoff2013} remain
possible.
\item \emph{Coding on bounded-depth side information, or depth-robust dependencies.} For an oblivious machine, Theorem~D
(\cref{thm:D}) leaves two ways out. Either some genie of $\eps L$ bits cuts the time--space dependency DAG to bounded
depth, and then the side information of the cones admits index codes of broadcast length $\eps L$, defeating the crux;
or no such genie exists, and the walks of the work heads are depth-robust. Home-based heads always admit such a genie
(Theorem~B, \cref{thm:B}), so for them only the first way is open (Corollary~B, \cref{cor:B}), and by the trade-off
(\cref{prop:tradeoff}) the depth is at least about $\beta/\eps$.
\item \emph{Codes that beat the pairing, of a particular kind.} If, as in the hardest instances we know, the relevant side
information is sparse, symmetric and locally tree-like, a construction needs storage codes of rate at least
$\tfrac12+\delta$ on it, for a fixed $\delta>0$ (a heuristic reading of \cref{prob:undirected-crux}). The classes of codes
that our results exclude are listed in \cref{sec:con-fail}.
Over a finite field, with a fixed number of symbols per vertex, a gain must come from $p$-torsion (\cref{cor:torsion}).
The binary code of rate $5/8$ on the Clebsch graph (\cref{cor:H4}) is not a torsion effect: it is the reduction of a rational
eigenvalue code with coefficients~$1$ (\cref{rem:ana-clebsch}), a mechanism that \cref{cor:torsion} rules out on locally
bipartite families of bounded degree. In odd
characteristic at most one of the two signed versions $I\pm W$ of a code can gain (\cref{prop:oddchar}).
\item \emph{If the dependencies are depth-robust: the deep-nesting template.} The only candidate we know for depth-robust walk pairs is
the deep-nesting pair (\cref{def:dnp}). Proof sketches indicate that per-head band rules fail already on one deeply nested walk
(\cref{rem:dep-N,rem:dep-Nprime}), although one such walk alone is easy to cut (\cref{lem:dep-visit}); and with one level
of nesting, labels exact on one head pay a price per hop on the cover independent of the depth (Theorem~E$_1$,
\cref{thm:E1}). Whether the pairs are depth-robust as $m$ grows is open (\cref{conj:Aprime}).
\end{enumerate}

\subsection{An architecture compatible with these conditions}

We describe, as a design and not as a result, how the first paper's geometric recursion~\cite[\S13]{FirstPaper2026} would
have to look in the light of (C1)--(C4). The recursion runs $\Theta(\log d)$ levels of geometrically decreasing volume,
each with $O(1)$ passes of an encoder and a decoder. Condition~(C1) forces the heads of each level to follow curved
rulers, for example near-identity Möbius maps realizing a monotone expander, and the second paper excludes the affine
variants~\cite{SecondPaper2026}. The encoder of each level must then save a constant fraction of the data by coding across
the expander's relations. By (C3), the code it uses must be nonlinear, or linear over a finite field,
with a $p$-torsion gain if it has a fixed number of symbols per vertex, and
a scalar linear code needs degree at least~$4$; its side-information graph must not map to any good target. Codes that beat $\tfrac12$ exist on single graphs with short
odd cycles, such as the Clebsch graph (girth~$4$) and, outside characteristic~$2$, the Hoffman--Singleton graph
(girth~$5$, \cref{prop:HS}). We do not know any family of bounded degree and growing odd girth that carries one; such a
family would answer \cref{prob:undirected-crux} in the negative. For an oblivious design, Corollary~B applies if its
work heads are all home-based, and the design then stands or falls with the crux; otherwise, if the crux holds, its
dependency structure must be depth-robust, and the deep-nesting pair is the template.

\subsection{Mechanisms that fail}\label{sec:con-fail}

The following classes of constructions fail for reasons proved here, in addition to those listed in the first two papers.
\begin{itemize}[leftmargin=2em]
\item \emph{Home-based motion of all work heads with legs of boundedly many constant-rate runs}, obliviously or
adaptively, with any number of legs at any scales (Theorem~C).
\item \emph{One expanding tape, affine returns, a common change of clock, or fixed changes of scale}: no single work tape
can carry the expansion, returns given by countably many affine maps in a common clock cannot carry it either, and
neither a common reparametrization of time nor a fixed reparametrization of each tape adds anything
(\cref{cor:mot-one-tape,thm:Cr,cor:mot-clock}).
\item \emph{Stroke trees for all work heads}, the recursive schedules of \cref{def:mot-stroke}: their limits are affine
wherever the heads move (\cref{prop:mot-strokes}).
\item \emph{Linear codes with bounded real or complex coefficients}, scalar or with a bounded number of symbols per
vertex, on locally bipartite side information of bounded degree (Theorem~A and \cref{prop:ana-A-ext}),
\emph{scalar linear codes of degree~$3$} on triangle-free side information (Theorem~S), and \emph{all codes} on
side information that is $C_5$-, $K_{14/5}$- or Petersen-colourable (\cref{cor:H2}).
\end{itemize}
